% 20260610_Assingment_ComputationalMathematics_A
\documentclass[dvipdfmx]{jsreport}
\usepackage{soupreamble}

\title{計算数学A 第3回レポート課題}
\author{M1 6012 川村聡一郎}
\date{}


\begin{document}
\maketitle
	グラフ$G= (V, E)$のサイズを$|V|+ |E|$で定義する。
\begin{tcolorbox}
\begin{enumerate}
	\item $n$個の閉区間に対する区間スケジュールを最大独立集合に変換したとき、変換後のグラフのサイズは最大どの程度になるか(オーダー表記でよい。)
	% Hint：|V|はnになる。|E|はどうなるか？例を挙げて示せ。
	\item 妻帯独立集合を$O(|V|)$で解くアルゴリズムが存在したとしても、区間スケジュールを$O(n)$で解けるとは言えない。理由をつけて答えよ。
	% Hint：区間スケジュールからソートへの変換にO(n)かかるならば、ソートがより小さいO(n')になったとしても区間スケジュールがO(n')にはならない。
\end{enumerate}
\end{tcolorbox}

\begin{souanswer} %答案はここに書く
	答案はここに書く
\end{souanswer}












\end{document}
